\documentclass{situationist-preprint}

\title{A Cosmological Argument From Well-Founded Dependence}
\msid{REL-01-d}
\msfamily{REL-01}
\msdiscipline{Philosophy of Religion}
\mstype{Alternative proof}
\mskeywords{cosmological argument; existential dependence; necessary being; well-foundedness; possibility premise}

\begin{document}
\maketitle

\begin{abstract}
The modal argument of \emph{On the Existence of God} (REL-01-a) rests on one substantive premise, the possibility that a maximally great being exists, and on the S5 conditions on accessibility. We give a route to the same conclusion that uses neither. Its premises are that no plurality of contingent beings is existentially self-sufficient, and that a necessary being on which contingent beings depend has maximal excellence essentially. The second premise permits infinite chains of dependence and is therefore consistent with the denial of a fundamental level. The argument establishes that a maximally great being exists, and it leaves the causal role of that being undetermined.
\end{abstract}

\section{The premise avoided}

The main argument of REL-01-a derives the necessary existence of a maximally great being (Theorem~3.4 of REL-01-a) from the possibility premise (Assumption~3.1) by way of the S5 conditions (Assumption~2.1), which yield Lemma~3.3, the principle that whatever is possibly necessary is necessary. REL-01-a observes that the same framework validates the converse argument from the possibility that no maximally great being exists (Proposition~4.1), so that the main argument stands or falls with a choice between two possibility premises. The argument below avoids the possibility premise altogether, and with it the symmetry: none of its premises has a converse that the framework treats on a par with it. It also makes no use of the symmetry or transitivity of accessibility. It belongs to the cosmological family \citep{Rowe1975, Pruss2006} and develops the convergence that the catalogue conclusion of REL-01 records.

\section{The argument}

We keep the worlds, the actual world $@$, the existence predicate and the definitions of maximal excellence and maximal greatness of REL-01-a (Definitions~2.2 to~2.4). A being is \emph{necessary} if it exists at every world accessible from $@$, and \emph{contingent} otherwise. We write $D(x,y)$ for: $x$ depends for its existence on $y$, in the sense of existential dependence \citep{Fine2010}. Dependence is not causation: it is the relation in which a set stands to its members, or a boundary to the region it bounds, and it holds without anything happening.

\begin{assumption}[Factivity]\label{a:fact}
If $D(x,y)$ holds at a world, both $x$ and $y$ exist at that world.
\end{assumption}

\begin{assumption}[Well-founded dependence]\label{a:wf}
For every non-empty plurality $X$ of contingent beings existing at $@$, some member of $X$ depends at $@$ on something that is not a member of $X$.
\end{assumption}

\begin{assumption}[Identification]\label{a:ident}
A necessary being on which some contingent being depends is maximally excellent at every world at which it exists.
\end{assumption}

\begin{assumption}[Contingency]\label{a:cont}
Some contingent being exists at $@$.
\end{assumption}

\begin{proposition}\label{p:necessary}
Some necessary being exists at $@$, and some contingent being depends on it.
\end{proposition}

\begin{proof}
Let $X$ be the plurality of all contingent beings existing at $@$; it is non-empty by Assumption~\ref{a:cont}. By Assumption~\ref{a:wf}, some member $c$ of $X$ depends at $@$ on some $y$ not in $X$. By Assumption~\ref{a:fact}, $y$ exists at $@$. Since $y$ exists at $@$ and is not among the contingent beings existing there, $y$ is necessary.
\end{proof}

\begin{theorem}\label{t:main}
Some being is maximally great at $@$, and therefore exists and is maximally excellent at $@$.
\end{theorem}

\begin{proof}
Let $y$ be as in Proposition~\ref{p:necessary}. Since $y$ is necessary, it exists at every world accessible from $@$; by Assumption~\ref{a:ident}, it is maximally excellent at each of them. By Definition~2.4 of REL-01-a, $y$ is maximally great at $@$. Reflexivity of accessibility gives existence and maximal excellence at $@$.
\end{proof}

Applied to a single contingent being, Assumption~\ref{a:wf} yields the familiar premise that every contingent being depends on something; the proof needs only its application to the plurality of all contingent beings. Reflexivity of accessibility is used only in the last step. Neither Assumption~3.1 of REL-01-a nor Lemma~3.3 appears.

\section{Objections}

\begin{objection}
Assumption~\ref{a:wf} is the claim that there is a fundamental level, which MET-10 denies.
\end{objection}

\begin{reply}
It is not. Assumption~\ref{a:wf} does not require that any member of $X$ depend on \emph{nothing} in $X$; it requires only that the plurality not be closed under dependence. An infinite descending chain of contingent beings, each depending on the next, satisfies it, provided that some member also depends on something outside the chain. The thesis of MET-10, that every level is grounded in a lower one without vicious regress, concerns the order among levels of contingent reality, and is consistent with the whole of that order depending on a being outside it \citep{Bliss2013}. What Assumption~\ref{a:wf} excludes is a plurality of contingent beings whose existence is accounted for entirely from within, which is the thesis that \citet[§§36--38]{Leibniz1714} took to leave the series as a whole unexplained.
\end{reply}

\begin{objection}
Assumption~\ref{a:ident} is as strong as the possibility premise it replaces.
\end{objection}

\begin{reply}
It is a different premise, and its support is independent. It says nothing about which properties are possibly instantiated; it ties maximal excellence to a role, that of a necessary being on which the existence of contingent beings depends. Its support is the traditional reasoning that a being on which contingent existence depends, and which exists without depending on anything contingent, is subject to no limitation imposed from outside \citep[ch.~2]{Rowe1975}. The premise belongs to Volume~II of REL-01, which REL-01-a describes as showing that considerations of metaphysical dependence converge on the object of the main argument. What it lacks is a converse of equal standing: the denial of Assumption~\ref{a:ident} is not supported by any consideration that the framework treats as symmetric with it.
\end{reply}

\section{What the argument leaves open}

The argument shows that a maximally great being exists and that some contingent being depends on it. It does not show that this being causes anything. Existential dependence holds without causal action, and Theorem~\ref{t:main} is consistent with models in which the dependee performs no act at all. The cosmological route therefore converges on the object of the modal argument without settling its causal role, as the catalogue conclusion of REL-01 states; REL-01-e draws the consequences of that limit for the main result.

\begin{thebibliography}{99}\small
\bibitem[Bliss(2013)]{Bliss2013} Bliss, R.~L. (2013). Viciousness and the structure of reality. \emph{Philosophical Studies}, 166(2), 399--418.
\bibitem[Fine(2010)]{Fine2010} Fine, K. (2010). Some puzzles of ground. \emph{Notre Dame Journal of Formal Logic}, 51(1), 97--118.
\bibitem[Leibniz(1714)]{Leibniz1714} Leibniz, G.~W. (1714). \emph{Monadologie}.
\bibitem[Pruss(2006)]{Pruss2006} Pruss, A.~R. (2006). \emph{The Principle of Sufficient Reason: A Reassessment}. Cambridge University Press.
\bibitem[Rowe(1975)]{Rowe1975} Rowe, W.~L. (1975). \emph{The Cosmological Argument}. Princeton University Press.
\end{thebibliography}

\end{document}
